% Generado por sistema/backend ReporteTex.java -- no editar a mano.
% Datos: Base de datos: 148 importado (titulados) / Base de datos: 53 importado (empleadores)
\textit{Paso 1. Variables.} Sobre los titulados que hicieron formación complementaria, se cruzan dos variables cualitativas -- $X$ = nivel del programa cursado, $Y$ = fuente de financiamiento -- con una variable cuantitativa -- $Z$ = años de vida profesional --, siguiendo la tabulación cruzada multivariada de la asignatura: frecuencia ($n$) y promedio de $Z$ ($\bar{Z}$) en cada combinación $X$--$Y$ que aparece en los datos.

\textit{Paso 2. Tabulación cruzada de frecuencias y antigüedad promedio.}

\begin{mitabla}{Nivel del programa ($X$) y fuente de financiamiento ($Y$), con antigüedad promedio ($n=96$)}
\begin{tabular}{llcc}
\toprule
$X$ (nivel) & $Y$ (financiamiento) & $n$ & $\bar{Z}$ (años) \\
\midrule
Diplomado & Recursos propios & 61 & 6.85 \\
Maestría & Recursos propios & 22 & 11.95 \\
Maestría & Beca parcial & 3 & 10.00 \\
Especialidad & Recursos propios & 3 & 12.33 \\
Doctorado & Recursos propios & 2 & 16.00 \\
Diplomado & Beca parcial & 2 & 6.50 \\
Maestría & Beca total & 2 & 25.50 \\
Diplomado & Beca total & 1 & 13.00 \\
\midrule
\multicolumn{2}{l}{Total / Promedio general} & 96 & 8.93 \\
\bottomrule
\end{tabular}
\end{mitabla}

\textit{Paso 3. Hallazgos por categoría de $X$.} Para cada nivel de programa, la fuente de financiamiento más frecuente y si su antigüedad promedio queda por encima o por debajo del promedio general (8.93 años):

\begin{itemize}
  \item \textbf{Doctorado} ($n=2$): antigüedad promedio de 16.00 años, por encima del promedio general. La fuente de financiamiento más frecuente es recursos propios (100.0\% de este grupo).
  \item \textbf{Maestría} ($n=27$): antigüedad promedio de 12.74 años, por encima del promedio general. La fuente de financiamiento más frecuente es recursos propios (81.5\% de este grupo).
  \item \textbf{Especialidad} ($n=3$): antigüedad promedio de 12.33 años, por encima del promedio general. La fuente de financiamiento más frecuente es recursos propios (100.0\% de este grupo).
  \item \textbf{Diplomado} ($n=64$): antigüedad promedio de 6.94 años, por debajo del promedio general. La fuente de financiamiento más frecuente es recursos propios (95.3\% de este grupo).
\end{itemize}

\textit{Interpretación.} A mayor nivel del programa cursado, mayor la antigüedad promedio de quienes lo cursaron: quienes ya acumularon más años de vida profesional son los que llegan a programas de nivel más alto (maestría o doctorado), mientras que los titulados más recientes se concentran en diplomados. Los recursos propios dominan el financiamiento en todos los niveles.

